%% kapa: exam=1 round=21 year=2010 period=2 form=A subject=appraisal_law subject_ko=부동산관계법규 seq=9 number=49
%% source: Q-Net 감정평가사 기출문제 (cst00309) artlSeq=5024181 file=extracted/1008/5024181/2010년 제21회 감정평가사 1차시험 기출문제/2교시.hwp
%% conversion: kordoc 4.24.51 -> tools/gen1.py (auto); review=proofread-gpt-sol; flags=none
\begin{question}{49}
국토의 계획 및 이용에 관한 법령상 개발행위허가에 관한 설명으로 옳지 않은 것은? (단, 조례는 고려하지 않음)
\begin{choices}
  \item 개발행위허가의 신청내용이 도시계획사업의 시행에 지장을 주는 때에는 개발행위허가를 할 수 없다.
  \item 토지의 일부를 공공용지 또는 공용지로 하기 위한 토지의 분할은 개발행위허가를 받을 필요가 없다.
  \item 개발행위허가권자는 개발행위허가 신청에 대하여 특별한 사유가 없으면 15일 이내에 허가 또는 불허가의 처분을 하여야 하는데, 이 15일에는 도시계획위원회의 심의나 관계행정기관의 장의 협의에 소요되는 기간은 포함되지 않는다.
  \item 개발행위허가권자는 개발행위를 허가하는 경우에 개발행위에 따른 기반시설의 설치, 환경오염방지, 경관, 조경 등에 관한 조치를 할 것을 조건으로 붙일 수 있으며, 이 때 미리 개발행위허가 신청자의 의견을 들을 필요는 없다.
  \item 개발행위허가권자는 공공시설의 귀속에 관한 사항이 포함된 개발행위허가를 하려면 미리 해당 공공시설이 속한 관리청의 의견을 들어야 한다.
\end{choices}
\end{question}
\answer{4}
